\paragraph{IV 21, subdivision 4: linear systems of divisors and invertible Modules (cote 25, sites 123--133).}
The heading on site 123 replaces a typed 7 by a handwritten 4. The opening calculation explicitly invokes subdivision 1, replaces a reference to subdivision 2 by \S20.3, and later cites Chapter III, \S2. It describes a divisor on the product with the dual projective bundle by an invertible Module \(N\) and a section \(\phi\), with fiberwise regularity and a degree-one requirement. The calculation expresses \(N\) through an invertible Module \(L\) on \(X\), identifies \(L\) by a direct image of \(N(-1)\), and uses projection formulas to pass to a homomorphism \(g^*E\longrightarrow L\), or equivalently \(u:E\longrightarrow g_*L\). Several overwritten projection indices remain uncertain in the transcription. Proposition 4.1 on site 124 organizes this construction as a correspondence with pairs \((L,u)\); its regularity test is stated after field extension, and the geometrically integral-fiber case is expressed by injectivity of the corresponding map into \(H^0(X_s,L_s)\). A marginal addition requires nonempty projective fibers. The author also requests an explicit reminder of the divisor obtained from the section of \(L\otimes\mathcal O_{\check P}(1)\). Corollary 4.2 begins on the same page with a finitely presented representing Module \(Q\). Its continuation on site 125 cites III 7.7.6 for the existence of \(Q\), uses Nakayama to pass from fiberwise surjectivity to a quotient \(Q\longrightarrow\check E\), and describes the remaining equivalence by multiplication with invertible sections. An added note records the role of \(f_*\mathcal O_X\simeq\mathcal O_S\) and of the corresponding units. The alternative description uses the divisor functor of subdivision 2 and the morphism \(\check P=P(\check E)\longrightarrow P(Q)\); site 126 then presents the associated classes through the Grassmannian of \(Q\). The manuscript's internal reference ``comme dans 5.1.'' in 4.2, the printed ``sur X'' for its earlier projective bundle, and its changes between \(f\) and \(g\) are retained in the text. See \href{https://grothendieck.umontpellier.fr/web/viewer.html?file=../25.pdf#page=124}{site 123}, \href{https://grothendieck.umontpellier.fr/web/viewer.html?file=../25.pdf#page=125}{site 124}, and \href{https://grothendieck.umontpellier.fr/web/viewer.html?file=../25.pdf#page=126}{sites 125--126}.

Proposition 4.3 begins on site 126 and treats the morphism from the parameter bundle to \(\mathrm{Div}_{X/S}\), followed by the action of automorphisms of that parameter bundle. The first part uses an associated prime cycle and its global functions; the second part is formulated with surjectivity of the structural map and includes a handwritten functorial addition. The author supplies a concrete warning on site 127: for an open subscheme of \(P^1_k\) omitting two distinct rational points, those two points give the same zero divisor in the displayed family. The proof first works over an algebraically closed field, passes to an associated cycle with the reduced induced structure, and compares sections differing by invertible functions. The lemma on site 128 chooses \(r+1\) points whose images form a projective basis. Its use on site 129 is explicitly tied to subdivision 1's description of automorphisms of a projective bundle. The general-base argument then passes through the noetherian case and local artinian quotients. On sites 130--131 the difference of two sections is expressed through tangent vectors and differentials, with separate descriptions for the divisor functor and for the morphism functor. On site 130, the written definition \(F_o=F\otimes_A k\) is followed by the phrase ``fibre réduite''; the argument also says \(h_o\) is surjective immediately after naming the preceding homomorphism \(h_o^*\). Site 131 gives the Hom expression the subscript \(\underline O_{X_o}\) and the target \(\underline O_{P_o}\). These literal choices remain in the text. Grothendieck says that this part should refer to the calculations of the planned paragraph on infinitesimal extensions, whose number is being changed from 25 toward 26. He also records the dependence of those calculations on paragraph 16 and explains the intended order of the references. Sites 131--132 describe the resulting interpretation of a linear system as a subfunctor of the divisor functor, the uniqueness of the parameter-bundle isomorphism, and the proposed use of faithfully flat descent in Chapter V with twisted projective bundles. These are the manuscript's own proof steps and plans; the guide supplies no missing tangent-space argument. See \href{https://grothendieck.umontpellier.fr/web/viewer.html?file=../25.pdf#page=128}{sites 127--129}, \href{https://grothendieck.umontpellier.fr/web/viewer.html?file=../25.pdf#page=131}{sites 130--131}, and \href{https://grothendieck.umontpellier.fr/web/viewer.html?file=../25.pdf#page=133}{site 132}.

The closing proposition on sites 132--133 is referred to in the surrounding discussion as 4.4, although its displayed number has been canceled. It returns to the set \(\mathrm{Ens}(D)\) of divisors obtained from rational points of the parameter bundle. The page replaces a description by rational functions with one involving a vector subspace of meromorphic functions and regularity of its nonzero members; the successive cancellations and the later occurrence of ``pseudo-'' remain visible. The last part compares the subspaces by multiplication with a regular function and records the remaining scalar ambiguity. Grothendieck leaves the proof as an exercise using 4.1 and proposes moving this proposition ahead of 4.3. He then discusses the additional algebraic-closure condition and its geometric-integrality case, before adding the convention that an ensemble of the form \(\mathrm{Ens}(D)\) may also be called a linear system of divisors. The nearby marginal discussion of complete linear systems on site 126 remains partly clipped and uncertain; it is not completed from this closing proposition. See \href{https://grothendieck.umontpellier.fr/web/viewer.html?file=../25.pdf#page=133}{site 132} and \href{https://grothendieck.umontpellier.fr/web/viewer.html?file=../25.pdf#page=134}{site 133}.
